\documentclass[10pt,a4paper]{amsart}
\usepackage{amsmath,amssymb,amsthm,graphicx}
\usepackage{verbatim}
\newenvironment{reference}{\comment}{\endcomment}
%\newenvironment{reference}{}{}
\newenvironment{slogan}{\comment}{\endcomment}
\newenvironment{history}{\comment}{\endcomment}

% For commutative diagrams we use Xy-pic
\usepackage[all]{xy}

% We use 2cell for 2-commutative diagrams.
\xyoption{2cell}
% Theorem environments.
%
\theoremstyle{plain}
\newtheorem{theorem}[subsection]{Theorem}
\newtheorem{proposition}[subsection]{Proposition}
\newtheorem{lemma}[subsection]{Lemma}

\theoremstyle{definition}
\newtheorem{definition}[subsection]{Definition}
\newtheorem{example}[subsection]{Example}
\newtheorem{exercise}[subsection]{Exercise}
\newtheorem{situation}[subsection]{Situation}

\theoremstyle{remark}
\newtheorem{remark}[subsection]{Remark}
\newtheorem{remarks}[subsection]{Remarks}

\numberwithin{equation}{subsection}

% Macros
%
\def\lim{\mathop{\mathrm{lim}}\nolimits}
\def\colim{\mathop{\mathrm{colim}}\nolimits}
\def\Spec{\mathop{\mathrm{Spec}}}
\def\Hom{\mathop{\mathrm{Hom}}\nolimits}
\def\Ext{\mathop{\mathrm{Ext}}\nolimits}
\def\SheafHom{\mathop{\mathcal{H}\!\mathit{om}}\nolimits}
\def\SheafExt{\mathop{\mathcal{E}\!\mathit{xt}}\nolimits}
\def\Sch{\mathit{Sch}}
\def\Mor{\mathop{\mathrm{Mor}}\nolimits}
\def\Ob{\mathop{\mathrm{Ob}}\nolimits}
\def\Sh{\mathop{\mathit{Sh}}\nolimits}
\def\NL{\mathop{N\!L}\nolimits}
\def\CH{\mathop{\mathrm{CH}}\nolimits}
\def\proetale{{pro\text{-}\acute{e}tale}}
\def\etale{{\acute{e}tale}}
\def\QCoh{\mathit{QCoh}}
\def\Ker{\mathop{\mathrm{Ker}}}
\def\Im{\mathop{\mathrm{Im}}}
\def\Coker{\mathop{\mathrm{Coker}}}
\def\Coim{\mathop{\mathrm{Coim}}}

% Boxtimes
%
\DeclareMathSymbol{\boxtimes}{\mathbin}{AMSa}{"02}

%
% Macros for moduli stacks/spaces
%
\def\QCohstack{\mathcal{QC}\!\mathit{oh}}
\def\Cohstack{\mathcal{C}\!\mathit{oh}}
\def\Spacesstack{\mathcal{S}\!\mathit{paces}}
\def\Quotfunctor{\mathrm{Quot}}
\def\Hilbfunctor{\mathrm{Hilb}}
\def\Curvesstack{\mathcal{C}\!\mathit{urves}}
\def\Polarizedstack{\mathcal{P}\!\mathit{olarized}}
\def\Complexesstack{\mathcal{C}\!\mathit{omplexes}}
% \Pic is the operator that assigns to X its picard group, usage \Pic(X)
% \Picardstack_{X/B} denotes the Picard stack of X over B
% \Picardfunctor_{X/B} denotes the Picard functor of X over B
\def\Pic{\mathop{\mathrm{Pic}}\nolimits}
\def\Picardstack{\mathcal{P}\!\mathit{ic}}
\def\Picardfunctor{\mathrm{Pic}}
\def\Deformationcategory{\mathcal{D}\!\mathit{ef}}
\usepackage{cleveref}
\crefname{originalequation}{original equation}{original equations}
\Crefname{originalequation}{Original equation}{Original equations}
\makeatletter
\newlabel{equation-plain-versus-derived-algebras}{{0BZ4}{source}{}{sourcepointer-1}{}}
\newlabel{equation-plain-versus-derived-algebras@cref}{{[originalequation][0][]0BZ4}{[1][0][]source}{}{}{}}
\makeatother
\crefname{originalreference}{original reference}{original references}
\Crefname{originalreference}{Original reference}{Original references}
\makeatletter
\newlabel{lemma-K-flat-resolution}{{06Y4}{source}{}{sourcepointer-2}{}}
\newlabel{lemma-K-flat-resolution@cref}{{[originalreference][0][]06Y4}{[1][0][]source}{}{}{}}
\makeatother
\crefname{originalreference}{original reference}{original references}
\Crefname{originalreference}{Original reference}{Original references}
\makeatletter
\newlabel{lemma-compose-tensor-functors}{{09S4}{source}{}{sourcepointer-3}{}}
\newlabel{lemma-compose-tensor-functors@cref}{{[originalreference][0][]09S4}{[1][0][]source}{}{}{}}
\makeatother
\crefname{originalreference}{original reference}{original references}
\Crefname{originalreference}{Original reference}{Original references}
\makeatletter
\newlabel{lemma-compose-tensor-functors-general}{{0BZ3}{source}{}{sourcepointer-4}{}}
\newlabel{lemma-compose-tensor-functors-general@cref}{{[originalreference][0][]0BZ3}{[1][0][]source}{}{}{}}
\makeatother
\crefname{originalreference}{original reference}{original references}
\Crefname{originalreference}{Original reference}{Original references}
\makeatletter
\newlabel{lemma-compose-tensor-functors-general-algebra}{{0BZ5}{source}{}{sourcepointer-5}{}}
\newlabel{lemma-compose-tensor-functors-general-algebra@cref}{{[originalreference][0][]0BZ5}{[1][0][]source}{}{}{}}
\makeatother
\crefname{originalreference}{original reference}{original references}
\Crefname{originalreference}{Original reference}{Original references}
\makeatletter
\newlabel{lemma-qis-equivalence}{{09S6}{source}{}{sourcepointer-6}{}}
\newlabel{lemma-qis-equivalence@cref}{{[originalreference][0][]09S6}{[1][0][]source}{}{}{}}
\makeatother
\crefname{originallemma}{original lemma}{original lemmas}
\Crefname{originallemma}{Original lemma}{Original lemmas}
\makeatletter
\newlabel{more-algebra-lemma-K-flat-quasi-isomorphism}{{06Y0}{source}{}{sourcepointer-7}{}}
\newlabel{more-algebra-lemma-K-flat-quasi-isomorphism@cref}{{[originallemma][0][]06Y0}{[1][0][]source}{}{}{}}
\makeatother
\title{Composition of derived tensor functors through an intermediate differential graded algebra is a derived tensor with a bimodule}
\author{The Stacks Project Authors}
\date{stacks-project:dga.tex:local-snapshot-20261012; GNU FDL 1.2 or later}
\begin{document}
\maketitle
\noindent\textit{Editorial note.} All delivered mathematical text is the author's own, copied from the pinned Stacks Project chapter file sources/stacks-dga/pinned-original-source.tex; no statement or proof step was rewritten, reordered or removed.\par\medskip
\noindent\textit{Editorial note.} The counted proof is the author's complete argument: the K-flat replacement with its equivalence and restriction quasi-inverse, the three-composition reduction, the dispatch of the two compositions, the property-(P) resolution P = L tensor\_R B with the rewriting of the comparison map, and the final reduction through more-algebra-lemma-K-flat-quasi-isomorphism to the standing hypothesis are all retained verbatim.\par\medskip
\noindent\textit{Editorial note.} The label lemma-compose-tensor-functors-tor names the counted statement, and the proof environment immediately following it in the source is delivered verbatim as its proof.\par\medskip
\noindent\textit{Editorial note.} The counted statement is carried by the author's own lemma environment in the chapter source.\par\medskip
\noindent\textit{Allowed standard prerequisites (editorial).}\begin{itemize}
\item K-flat replacement of the intermediate algebra. Same chapter, Lemmas lemma-K-flat-resolution and lemma-qis-equivalence. Lets the intermediate algebra (B, d) be replaced by a K-flat quasi-isomorphic (B', d) whose derived category is equivalent to D(B, d) with restriction as quasi-inverse.
\item Reduction lemmas for plain versus derived tensor. Same chapter, Lemmas lemma-compose-tensor-functors, lemma-compose-tensor-functors-general-algebra; More on Algebra, Lemma more-algebra-lemma-K-flat-quasi-isomorphism. Handle the two reduced compositions, identify the target map for the K-flat case, and finally replace a quasi-isomorphism tensored with a K-flat complex.
\end{itemize}\medskip
\ifdefined\hypertarget\hypertarget{sourcepointer-1}{}\fi
\noindent\textit{Source reference.} Original equation 0BZ4: sources/stacks-dga/pinned-original-source.tex, source lines 5941--5941. Prerequisite pointer invoked at proof line(s) 7255; it names a standard fact of the same chapter that the counted argument uses, and every step of the counted proof built on it is retained verbatim in the delivered body.\par\medskip
\ifdefined\hypertarget\hypertarget{sourcepointer-2}{}\fi
\noindent\textit{Source reference.} Original reference 06Y4: sources/stacks-dga/pinned-original-source.tex, source lines 7157--7157. Prerequisite pointer invoked at proof line(s) 7230; it names a standard fact of the same chapter that the counted argument uses, and every step of the counted proof built on it is retained verbatim in the delivered body.\par\medskip
\ifdefined\hypertarget\hypertarget{sourcepointer-3}{}\fi
\noindent\textit{Source reference.} Original reference 09S4: sources/stacks-dga/pinned-original-source.tex, source lines 6067--6067. Prerequisite pointer invoked at proof line(s) 7245; it names a standard fact of the same chapter that the counted argument uses, and every step of the counted proof built on it is retained verbatim in the delivered body.\par\medskip
\ifdefined\hypertarget\hypertarget{sourcepointer-4}{}\fi
\noindent\textit{Source reference.} Original reference 0BZ3: sources/stacks-dga/pinned-original-source.tex, source lines 5867--5867. Prerequisite pointer invoked at proof line(s) 7249; it names a standard fact of the same chapter that the counted argument uses, and every step of the counted proof built on it is retained verbatim in the delivered body.\par\medskip
\ifdefined\hypertarget\hypertarget{sourcepointer-5}{}\fi
\noindent\textit{Source reference.} Original reference 0BZ5: sources/stacks-dga/pinned-original-source.tex, source lines 5959--5959. Prerequisite pointer invoked at proof line(s) 7257; it names a standard fact of the same chapter that the counted argument uses, and every step of the counted proof built on it is retained verbatim in the delivered body.\par\medskip
\ifdefined\hypertarget\hypertarget{sourcepointer-6}{}\fi
\noindent\textit{Source reference.} Original reference 09S6: sources/stacks-dga/pinned-original-source.tex, source lines 6725--6725. Prerequisite pointer invoked at proof line(s) 7233; it names a standard fact of the same chapter that the counted argument uses, and every step of the counted proof built on it is retained verbatim in the delivered body.\par\medskip
\ifdefined\hypertarget\hypertarget{sourcepointer-7}{}\fi
\noindent\textit{Source reference.} Original lemma 06Y0: sources/stacks-chow/tags.txt, source lines 8869--8869. Prerequisite pointer invoked at proof line(s) 7274; it names a standard fact of a sibling chapter that the counted argument uses, and every step of the counted proof built on it is retained verbatim in the delivered body.\par\medskip
\begin{lemma}
\label{lemma-compose-tensor-functors-tor}
Let $R$ be a ring. Let $(A, \text{d})$, $(B, \text{d})$, and
$(C, \text{d})$ be differential graded $R$-algebras. Assume
$A \otimes_R C$ represents $A \otimes^\mathbf{L}_R C$ in $D(R)$.
Let $N$ be a differential graded $(A, B)$-bimodule.
Let $N'$ be a differential graded $(B, C)$-bimodule.
Then the composition
$$
\xymatrix{
D(A, \text{d}) \ar[rr]^{- \otimes_A^\mathbf{L} N} & &
D(B, \text{d}) \ar[rr]^{- \otimes_B^\mathbf{L} N'} & &
D(C, \text{d})
}
$$
is isomorphic to $- \otimes_A^\mathbf{L} N''$ for some differential graded
$(A, C)$-bimodule $N''$.
\end{lemma}

\begin{proof}
Using Lemma \ref{lemma-K-flat-resolution}
we choose a quasi-isomorphism $(B', \text{d}) \to (B, \text{d})$
with $B'$ K-flat as a complex of $R$-modules.
By Lemma \ref{lemma-qis-equivalence}
the functor $-\otimes^\mathbf{L}_{B'} B : D(B', \text{d}) \to D(B, \text{d})$
is an equivalence with quasi-inverse given by restriction.
Note that restriction is canonically isomorphic to the functor
$- \otimes^\mathbf{L}_B B : D(B, \text{d}) \to D(B', \text{d})$
where $B$ is viewed as a $(B, B')$-bimodule.
Thus it suffices to prove the lemma for the compositions
$$
D(A) \to D(B) \to D(B'),\quad
D(B') \to D(B) \to D(C),\quad
D(A) \to D(B') \to D(C).
$$
The first one is Lemma \ref{lemma-compose-tensor-functors}
because $B'$ is K-flat as a complex of $R$-modules.
The second one is true because
$B \otimes_B^\mathbf{L} N' = N' = B \otimes_B N'$
and hence Lemma \ref{lemma-compose-tensor-functors-general} applies.
Thus we reduce to the case where $B$ is K-flat as a complex
of $R$-modules.

\medskip\noindent
Assume $B$ is K-flat as a complex of $R$-modules. It suffices to
show that (\ref{equation-plain-versus-derived-algebras}) is an
isomorphism, see
Lemma \ref{lemma-compose-tensor-functors-general-algebra}.
Choose a quasi-isomorphism $L \to A$ where $L$ is a differential
graded $R$-module which has property (P). Then it is clear that
$P = L \otimes_R B$ has property (P) as a differential graded $B$-module.
Hence we have to show that $P \to A \otimes_R B$
induces a quasi-isomorphism
$$
P \otimes_B (B \otimes_R C)
\longrightarrow
(A \otimes_R B) \otimes_B (B \otimes_R C)
$$
We can rewrite this as
$$
P \otimes_R B \otimes_R C \longrightarrow A \otimes_R B \otimes_R C
$$
Since $B$ is K-flat as a complex of $R$-modules, it
follows from
More on Algebra, Lemma \ref{more-algebra-lemma-K-flat-quasi-isomorphism}
that it is enough
to show that
$$
P \otimes_R C \to A \otimes_R C
$$
is a quasi-isomorphism, which is exactly our assumption.
\end{proof}

\end{document}
